%%%%%%%%%%%%%%%%%%%%%%%%%%%%%%%%%%%%%%%%%%%%%%%
%
% Template per Elaborato di Laurea
% DISI - Dipartimento di Ingegneria e Scienza dell’Informazione
%
% update 2015-09-10
%
% Per la generazione corretta del 
% pdflatex nome_file.tex
% bibtex nome_file.aux
% pdflatex nome_file.tex
% pdflatex nome_file.tex
%
%%%%%%%%%%%%%%%%%%%%%%%%%%%%%%%%%%%%%%%%%%%%%%%

% formato FRONTE RETRO
\documentclass[epsfig,a4paper,11pt,titlepage,twoside,openany]{book}
\usepackage{epsfig}
\usepackage{plain}
\usepackage{setspace}
\usepackage[paperheight=29.7cm,paperwidth=21cm,outer=1.5cm,inner=2.5cm,top=2cm,bottom=2cm]{geometry} % per definizione layout
\usepackage{titlesec} % per formato custom dei titoli dei capitoli

%%%%%%%%%%%%%%
% supporto lettere accentate
%
%\usepackage[latin1]{inputenc} % per Windows;
\usepackage[utf8x]{inputenc} % per Linux (richiede il pacchetto unicode);
%\usepackage[applemac]{inputenc} % per Mac.

\singlespacing

\usepackage[english]{babel}

\begin{document}

  % nessuna numerazione
  \pagenumbering{gobble} 
  \pagestyle{plain}

\thispagestyle{empty}

\begin{center}
  \begin{figure}[h!]
    \centerline{\psfig{file=marchio_unitrento_colore_it_202002.eps,width=0.6\textwidth}}
  \end{figure}

  \vspace{2 cm} 

  \LARGE{Department of Information Engineering and Computer Science\\}

  \vspace{1 cm} 
  \Large{Bachelor degree in\\
  Computer Science}

  \vspace{2 cm} 
  \Large\textsc{final dissertation\\} 
  \vspace{1 cm} 
  \Huge\textsc{\pyvmt{}\\}
  \Large{\it{A \python{} Library to manage \VMT{} Problems}}


  \vspace{2 cm} 
  \begin{tabular*}{\textwidth}{ c @{\extracolsep{\fill}} c }
  \Large{Supervisor} & \Large{Student}\\
  \Large{Sebastiani Roberto}& \Large{Girardi Alessandro}\\
  \\
  \Large{Co-supervisor} & \\
  \Large{Griggio Alberto}& \\
  \end{tabular*}

  \vspace{2 cm} 

  \Large{Academic year 2021/2022}
  
\end{center}



  \clearpage
 
%%%%%%%%%%%%%%%%%%%%%%%%%%%%%%%%%%%%%%%%%%%%%%%%%%%%%%%%%%%%%%%%%%%%%%%%%%
%%%%%%%%%%%%%%%%%%%%%%%%%%%%%%%%%%%%%%%%%%%%%%%%%%%%%%%%%%%%%%%%%%%%%%%%%%
%% Nota
%%%%%%%%%%%%%%%%%%%%%%%%%%%%%%%%%%%%%%%%%%%%%%%%%%%%%%%%%%%%%%%%%%%%%%%%%%
%% Sezione Ringraziamenti opzionale
%%%%%%%%%%%%%%%%%%%%%%%%%%%%%%%%%%%%%%%%%%%%%%%%%%%%%%%%%%%%%%%%%%%%%%%%%%
%%%%%%%%%%%%%%%%%%%%%%%%%%%%%%%%%%%%%%%%%%%%%%%%%%%%%%%%%%%%%%%%%%%%%%%%%%
  \thispagestyle{empty}

\begin{center}
  {\bf \Huge Acknowledgements}
\end{center}

\vspace{4cm}


\emph{
  ...thanks to...
}

  \clearpage
  \pagestyle{plain} % nessuna intestazione e pie pagina con numero al centro

  
  % inizio numerazione pagine in numeri arabi
  \mainmatter

%%%%%%%%%%%%%%%%%%%%%%%%%%%%%%%%%%%%%%%%%%%%%%%%%%%%%%%%%%%%%%%%%%%%%%%%%%
%%%%%%%%%%%%%%%%%%%%%%%%%%%%%%%%%%%%%%%%%%%%%%%%%%%%%%%%%%%%%%%%%%%%%%%%%%
%% Nota
%%%%%%%%%%%%%%%%%%%%%%%%%%%%%%%%%%%%%%%%%%%%%%%%%%%%%%%%%%%%%%%%%%%%%%%%%%
%% Si ricorda che il numero massimo di facciate e' 30.
%% Nel conteggio delle facciate sono incluse 
%%   indice
%%   sommario
%%   capitoli
%% Dal conteggio delle facciate sono escluse
%%   frontespizio
%%   ringraziamenti
%%   allegati    
%%%%%%%%%%%%%%%%%%%%%%%%%%%%%%%%%%%%%%%%%%%%%%%%%%%%%%%%%%%%%%%%%%%%%%%%%%
%%%%%%%%%%%%%%%%%%%%%%%%%%%%%%%%%%%%%%%%%%%%%%%%%%%%%%%%%%%%%%%%%%%%%%%%%%

    % indice
    \tableofcontents
    \clearpage
    
    
          
    % gruppo per definizone di successione capitoli senza interruzione di pagina
    \begingroup
      % nessuna interruzione di pagina tra capitoli
      % ridefinizione dei comandi di clear page
      \renewcommand{\cleardoublepage}{} 
      \renewcommand{\clearpage}{} 
      % redefinizione del formato del titolo del capitolo
      % da formato
      %   Capitolo X
      %   Titolo capitolo
      % a formato
      %   X   Titolo capitolo
      
      \titleformat{\chapter}
        {\normalfont\Huge\bfseries}{\thechapter}{1em}{}
        
      \titlespacing*{\chapter}{0pt}{0.59in}{0.02in}
      \titlespacing*{\section}{0pt}{0.20in}{0.02in}
      \titlespacing*{\subsection}{0pt}{0.10in}{0.02in}
      
      % sommario
      \chapter*{Abstract}
\label{abstract}

\addcontentsline{toc}{chapter}{Abstract}

\paragraph*{Context and motivation}
PyVmt was created by me as part of a formative internship at Fondazione Bruno Kessler. It's an extension of an existing project which was also created in the foundation, pySMT, which offers a solver-agnostic formula manipulation interface.
Several factors motivate the creation for an extension:
\begin{itemize}
  \item the creation of specific tools and functions which are commonly used in model verification
  \item the extension of the existing SMT-LIB parser to extract the information relating to the model and its properties written in VMT-LIB format
  \item the addition of support for model checkers to offer an experience similar to pySMT while dealing with verification problems
\end{itemize}

\paragraph*{What can PyVmt do?}
PyVmt is a python library that enables the user to manage VMT models by reading and writing code in the VMT-LIB format, modifying models by adding variables, constraints, properties, or with some higher level function such as renaming and synchronous composition. It offers several model checkers as backend for verification, making it simpler for the user to choose several of them and compare their performance.

\paragraph*{Techniques used}
PyVmt is written in Python, it was developed using a continuous integration approach, utilizing Github actions as a tool to verify non-regression after each update.
Tests are run using the pytest test suite and are mostly composed of unit tests which aim to cover most of the library's source code.
Most of the documentation is present in the source code as doc-strings so that they can easily be looked up while using an IDE. An external reference is also available and can be generated using sphinx.
After all the tests pass, the new version of the library is automatically uploaded as a package to PyPi so that any user can use it in their own project.

\paragraph*{Results reached}
To verify the functionality of PyVmt, I collected several benchmarks and ran them against the available solvers. The results from these experiments show a promising support from the backends to deal with most of the problems within their supported theories.

%%%%%%%%%%%%%%%%%%%%%%%%%%%%%%%%%%%%%%%%%%%%%%%%%%%%%%%%%%%%%%%%%%%%%%%%%%
%%%%%%%%%%%%%%%%%%%%%%%%%%%%%%%%%%%%%%%%%%%%%%%%%%%%%%%%%%%%%%%%%%%%%%%%%%
%% Nota
%%%%%%%%%%%%%%%%%%%%%%%%%%%%%%%%%%%%%%%%%%%%%%%%%%%%%%%%%%%%%%%%%%%%%%%%%%
%% Sommario e' un breve riassunto del lavoro svolto dove si descrive 
%% l’obiettivo, l’oggetto della tesi, le metodologie e 
%% le tecniche usate, i dati elaborati e la spiegazione delle conclusioni 
%% alle quali siete arrivati.
%% Il sommario dell’elaborato consiste al massimo di 3 pagine e deve contenere le seguenti informazioni: 
%%   contesto e motivazioni
%%   breve riassunto del problema affrontato
%%   tecniche utilizzate e/o sviluppate
%%   risultati raggiunti, sottolineando il contributo personale del laureando/a
%%%%%%%%%%%%%%%%%%%%%%%%%%%%%%%%%%%%%%%%%%%%%%%%%%%%%%%%%%%%%%%%%%%%%%%%%%
%%%%%%%%%%%%%%%%%%%%%%%%%%%%%%%%%%%%%%%%%%%%%%%%%%%%%%%%%%%%%%%%%%%%%%%%%%      
      
      %%%%%%%%%%%%%%%%%%%%%%%%%%%%%%%%
      % lista dei capitoli
      %
      % \input oppure \include
      %
      \chapter{Verification modulo theory}

\paragraph*{SMT}
Satisfiability modulo theory is a generalization of the SAT (boolean satisfiability) problem. It involves proving whether a formula can be satisfied by some assignment of the variables. The expressions, which can include lists, arrays, bit vectors, and strings are interpreted within a formal theory. There exists a standard language for defining SMT problems called SMT-LIB \cite{SMT-LIB} and it simplifies interfacing to the solvers, which are the programs developed to decide whether a formula is satisfiable.

\paragraph*{VMT}
Verification modulo theory adds a dynamic aspect to SMT, allowing the specification of transition systems. In VMT, transition systems are defined in a symbolic manner, they're composed of several elements:
\begin{itemize}
    \item State variables
    \item An init constraint
    \item A transition constraint
\end{itemize}
States are therefore not explicitly expressed, rather the assignments to the state variables may be used to encode them.
This has the advantage of allowing for transition systems with an extremely large number of states, and is frequently a more natural way to define their behavior.
The transition between states is regulated by a transition constraint, which is a function between current and next state variables and must be verified during a transition. The set of initial states is also not explicit, initial states must verify the init constraint.


\section{Importance of model checking}
The classical approach for finding and resolving bugs in software usually involves writing tests and code reviews done by other developers.
The safety of a system tested in this manner is not guaranteed since only a limited amount of scenarios can be explored, and subtle bugs may go unnoticed even after peer-review.
A mathematical approach is therefore necessary whenever such guarantees are required, for example when dealing with complex, safety-critical systems where bugs may have severe repercussions and are therefore not acceptable.

For these applications, software may be written utilizing a formal language, allowing for model checking to be applied in order to formally verify some properties corresponding to the specifications.

\section{Some applications}
Many different problems can be solved by using a verification tool, the following paragraphs include a few relevant ones.

\paragraph*{LTL} LTL or Linear Temporal Logic is used to define complex properties over execution sequences of the system.
It works with a fragment of first order logic and adds several operators (X, F, G, U) to introduce the concept of time to our properties.
The semantics of the operators is as follows:
\begin{itemize}
    \item \begin{math} X \phi \end{math}, in the next step \begin{math} \phi \end{math} is verified
    \item \begin{math} F \phi \end{math}, in some future step \begin{math} \phi \end{math} has to be verified
    \item \begin{math} G \phi \end{math}, \begin{math} \phi \end{math} is globally verified
    \item \begin{math} \phi\ R\ \psi \end{math}, \begin{math} \psi \end{math} is verified as long as \begin{math} \phi \end{math} is not verified
\end{itemize}

Combinations of these operators may be used to create classic property types or more complex properties.

\paragraph*{Safety} Safety properties are used to define that the system never reaches states in which a condition we want to avoid is verified, in LTL this can be written as \begin{math}G (\lnot bad)\end{math}.
This verifies that the states which the system can reach is limited to safe ones and that the bad behavior cannot be reached during the execution.

An example of a safety problem is ensuring that it's impossible for the door of a microwave to be open when the device is running:
\begin{math}
    G \lnot (open \land running)
\end{math}.

\paragraph*{Fairness} Fairness specifies that a set of states must be visited infinitely often during execution and is denoted in LTL as \begin{math} G(F(\phi))\end{math}.
Fairness can be used both as a property to be checked or as a constraint that defines how infinite traces of the system should look like.

A problem involving a fairness property is ensuring that a scheduler for an operating system eventually schedules each process, thus not leaving any process to wait infinitely often.

\paragraph*{Liveness} Liveness properties specify that after a while, a property holds forever. It's denoted in LTL as \begin{math} F(G(\phi)) \end{math}.
This kind of property is frequently explained with the notion that something good will eventually happen and relates to the concept of progress.
They're similar to fairness properties, in fact the previous formula is equivalent to \begin{math} \lnot G(F(\lnot \phi))\end{math} meaning that a liveness property may be proven invalid by a trace in which \begin{math} \phi \end{math} is negated infinitely often, which is a fair path for \begin{math}\lnot \phi\end{math}.

\paragraph*{Decision problems} Some complex logic problems involving decisions on an evolving system may be modelled as a transition system.
After that is done, a property can be added specifying that a solution to the original problem may not be found, i.e. \begin{math} G (\lnot solved) \end{math}.
A specialized program may be used to verify whether the property holds, if it doesn't (meaning that a solution to our problem exists), it may print the exact steps required to reach such solution.

A classic problem that can be solved in this manner is the ferryman problem, where a ferryman has to carry across a river several entities (e.g. a cabbage, a wolf, and a sheep) on a boat with a limited capacity, while never leaving certain pairs of entities unattended (e.g. the sheep with the cabbage or the sheep with the wolf).
Once modelled, the solution can be found as a counterexample to the property:
\begin{math}
    G (\exists (entity). \lnot isAcross(entity))
\end{math}.
The resulting counterexample displays which entities have to be carried across the river and back in each step.

\section{The VMT-LIB format}
VMT-LIB is an extension of the SMT2 format and provides a standard interface for defining transition systems by exploiting the possibility to attach annotations to terms and formulas, which can be used to specify the model behaviors. \cite{VMT-LIB}

It is designed with a limited amount of language constructs, aiming for simplicity for anyone interested in implementing the language within their VMT solver:
\begin{itemize}[noitemsep]
    \item \textbf{:next name} is used to specify state variables. To define a state variable, two variables are created, one representing the current state and one representing the next state. The next annotation is used to link the next state variable to the current state variable.
    \item \textbf{:init} is used to specify the formula used to constrain the initial states for the system.
    \item \textbf{:trans} is used to specify the formula used to constrain the transitions from each state to the next for the system.
    \item \textbf{:invar-property idx} is used to specify an invariant property to be verified, in LTL it's in the form \begin{math} G \phi \end{math} and requires a formula to always hold.  A unique non-negative integer specifies the index of the property.
    \item \textbf{:live-property idx} is used to specify a live property to be verified, in LTL it's in the form \begin{math} F\ G\ \phi\end{math} and requires a formula to hold infinitely often. A unique non-negative integer specifies the index of the property.
\end{itemize}
\cite{vmt-lib-paper} %...

% execution paths...


      \chapter{Conclusion}
\label{cha:conclusions}



      %\input{capitolo4}
      
      
    \endgroup


    % bibliografia in formato bibtex
    %
    % aggiunta del capitolo nell'indice
    \addcontentsline{toc}{chapter}{Bibliography}
    % stile con ordinamento alfabetico in funzione degli autori
    \bibliographystyle{plain}
    \bibliography{biblio}
%%%%%%%%%%%%%%%%%%%%%%%%%%%%%%%%%%%%%%%%%%%%%%%%%%%%%%%%%%%%%%%%%%%%%%%%%%
%%%%%%%%%%%%%%%%%%%%%%%%%%%%%%%%%%%%%%%%%%%%%%%%%%%%%%%%%%%%%%%%%%%%%%%%%%
%% Nota
%%%%%%%%%%%%%%%%%%%%%%%%%%%%%%%%%%%%%%%%%%%%%%%%%%%%%%%%%%%%%%%%%%%%%%%%%%
%% Nella bibliografia devono essere riportati tutte le fonti consultate 
%% per lo svolgimento della tesi. La bibliografia deve essere redatta 
%% in ordine alfabetico sul cognome del primo autore. 
%% 
%% La forma della citazione bibliografica va inserita secondo la fonte utilizzata:
%% 
%% LIBRI
%% Cognome e iniziale del nome autore/autori, la data di edizione, titolo, casa editrice, eventuale numero dell’edizione. 
%% 
%% ARTICOLI DI RIVISTA
%% Cognome e iniziale del nome autore/autori, titolo articolo, titolo rivista, volume, numero, numero di pagine.
%% 
%% ARTICOLI DI CONFERENZA
%% Cognome e iniziale del nome autore/autori (anno), titolo articolo, titolo conferenza, luogo della conferenza (città e paese), date della conferenza, numero di pagine. 
%% 
%% SITOGRAFIA
%% La sitografia contiene un elenco di indirizzi Web consultati e disposti in ordine alfabetico. 
%% E’ necessario:
%%   Copiare la URL (l’indirizzo web) specifica della pagina consultata
%%   Se disponibile, indicare il cognome e nome dell’autore, il titolo ed eventuale sottotitolo del testo
%%   Se disponibile, inserire la data di ultima consultazione della risorsa (gg/mm/aaaa).    
%%%%%%%%%%%%%%%%%%%%%%%%%%%%%%%%%%%%%%%%%%%%%%%%%%%%%%%%%%%%%%%%%%%%%%%%%%
%%%%%%%%%%%%%%%%%%%%%%%%%%%%%%%%%%%%%%%%%%%%%%%%%%%%%%%%%%%%%%%%%%%%%%%%%%
    

    % \titleformat{\chapter}
    %     {\normalfont\Huge\bfseries}{Allegato \thechapter}{1em}{}
    % sezione Allegati - opzionale
    % \appendix
    % \input{allegati}

\end{document}
